\section{Conclusion and Future Works}
In this paper, we studied the power of adaptivity for pure exploration in linear bandits. We established matching upper and lower bounds on the minimax sample complexity of non-adaptive fixed-design algorithms.We then identified structured action sets for which adaptivity yields at most a logarithmic improvement over our non-adaptive approach, and we constructed a structured action set for which adaptive sampling attains a polynomial improvement over any non-adaptive approach. To obtain the polynomial separation, we develop an $\ell_2$-norm estimation procedure that uses $\mathcal{O}\!\left(\frac{d\log(1/\delta)}{\varepsilon^2}\right)$ samples. As a consequence, we also obtain a polynomial separation between the sample complexity of identifying an $\varepsilon$-best arm and that of estimating the optimal value $\max_{x\in\mathcal{X}}\langle x,\theta\rangle$ when $\mathcal{X}$ is the unit $\ell_2$ ball.

Our work raises several open questions. Is there an intrinsic and interpretable property of $\mathcal{X}$ that determines whether adaptivity can yield at most logarithmic improvements or instead enables polynomial improvements over non-adaptive approaches? Is the minimax sample complexity of estimating $\max_{x\in\mathcal{X}}\langle x,\theta\rangle$ always $\mathcal{O}\!\left(\frac{d\log(1/\delta)}{\varepsilon^2}\right)$ for all action sets $\mathcal{X}\subset \mathbb{R}^d$? Can our value-estimation-based adaptive approach be extended to other structured sets to obtain polynomial improvements, and can our adaptive lower-bound techniques be generalized beyond the classes considered in this paper? 

{Finally, we hope these results help connect various lines of work: for experimental design, they suggest that pure exploration in linear bandits \emph{can} exhibit polynomial advantages from adaptivity under suitable action-set geometry; for reinforcement learning through a PAC-learning lens, they suggest that the structure of the feature mapping in linear function approximation (see, e.g., \cite{bradtke1996linear,jin2020provably}) may fundamentally determine when adaptive exploration across episodes is beneficial.
}
